\chapter{Glossary (50 Terms)}
\label{app:glossary}

This appendix defines key terms used throughout the thesis.

%==============================================================================
\section{Core Concepts}
%==============================================================================

\begin{table}[h]
\centering
\caption{Glossary (Selected Terms)}
\begin{tabular}{ll}
\toprule
Term & Definition \\
\midrule
AUROC & Area Under the ROC Curve; threshold-invariant ranking metric. \\
AUPRC & Area Under the Precision-Recall Curve. \\
Few-shot & Learning from few labeled examples per class. \\
FedAvg & Federated Averaging; aggregation of client model updates. \\
MIMIC-IV & Medical Information Mart for Intensive Care, version 4. \\
SOFA & Sequential Organ Failure Assessment score. \\
SHAP & SHapley Additive exPlanations; feature attribution method. \\
\midrule
\multicolumn{2}{l}{\textit{Representation}} \\
\midrule
Symptom Token & Discrete representation of lab/vital value (bin + intensity). \\
ETHOS & Encoder for Time-series and Heterogeneous Organ Signals. \\
Prototype & Class centroid in embedding space (prototypical networks). \\
\midrule
\multicolumn{2}{l}{\textit{Learning}} \\
\midrule
MAML & Model-Agnostic Meta-Learning. \\
Episodic & Training regime with support/query splits. \\
Focal Loss & Loss that down-weights easy examples. \\
\midrule
\multicolumn{2}{l}{\textit{Clinical}} \\
\midrule
Feasibility & Discharge to home/rehab vs.\ death/SNF. \\
ICU LOS & Intensive Care Unit Length of Stay. \\
ICD-10 & International Classification of Diseases, 10th revision. \\
\midrule
\multicolumn{2}{l}{\textit{Evaluation}} \\
\midrule
MCC & Matthews Correlation Coefficient. \\
F1-macro & Mean of per-class F1 scores. \\
Bootstrap CI & Confidence interval from resampling. \\
Wilcoxon & Non-parametric paired test. \\
\midrule
\multicolumn{2}{l}{\textit{Technical}} \\
\midrule
ItemID & MIMIC identifier for lab/vital type. \\
Stay ID & Unique ICU stay identifier. \\
Parquet & Columnar storage format. \\
\bottomrule
\end{tabular}
\label{tab:glossary_full}
\end{table}

%==============================================================================
\section{Extended Glossary (Full List)}
%==============================================================================

\begin{longtable}{p{2.2cm}p{10cm}}
\toprule
\textbf{Term} & \textbf{Definition} \\
\midrule
\endfirsthead

\toprule
\textbf{Term} & \textbf{Definition} \\
\midrule
\endhead

Adversarial robustness & Model resilience to perturbed inputs. \\
Attention mechanism & Weights over input tokens for context-dependent representation. \\
AUROC & Area Under the ROC Curve; probability that a positive ranks above a negative. \\
AUPRC & Area Under the Precision-Recall Curve; useful for imbalanced data. \\
Bootstrap CI & Confidence interval from resampling with replacement. \\
Calibration & Predicted probabilities match observed frequencies. \\
Class imbalance & Unequal class distribution (e.g., 42\% vs.\ 58\%). \\
Cross-validation & K-fold evaluation for robust performance estimates. \\
Differential privacy & Formal privacy guarantee for aggregated statistics. \\
Embedding & Dense vector representation of input (e.g., 128-D). \\
Episodic training & Meta-learning with support/query episode sampling. \\
ETHOS & Encoder for Time-series and Heterogeneous Organ Signals. \\
FedAvg & Federated Averaging; weighted average of client updates. \\
Federated learning & Training across distributed data without centralization. \\
Feasibility & Binary: discharge to home/rehab (1) vs.\ death/SNF (0). \\
Focal loss & Cross-entropy with $(1-p)^\gamma$ down-weighting of easy examples. \\
F1-macro & Mean of per-class F1; treats classes equally. \\
Gradient descent & Optimization by iteratively moving opposite to gradient. \\
Hyperparameter & Model/config parameter set before training (e.g., learning rate). \\
ICD-10 & International Classification of Diseases, 10th revision. \\
ICU LOS & Intensive Care Unit Length of Stay. \\
Imputation & Filling missing values (e.g., forward-fill, mean). \\
ItemID & MIMIC identifier for lab/vital measurement type. \\
K-shot & Number of support examples per class (e.g., K=5). \\
Label smoothing & Softening one-hot labels for regularization. \\
LightGBM & Gradient boosting with leaf-wise growth. \\
MAML & Model-Agnostic Meta-Learning; learn initialization for fast adaptation. \\
MCC & Matthews Correlation Coefficient; balanced for binary classification. \\
Meta-learning & Learning to learn; optimizing across tasks. \\
MIMIC-IV & Medical Information Mart for Intensive Care, version 4. \\
Normalization & Scaling inputs to standard range (e.g., [0,1]). \\
Outlier & Extreme value; we clip to clinical ranges. \\
Parquet & Columnar storage format for efficient I/O. \\
Precision & TP / (TP + FP); positive predictive value. \\
Prototype & Class centroid in embedding space (prototypical networks). \\
Recall & TP / (TP + FN); sensitivity. \\
Resampling & Oversampling minority or undersampling majority class. \\
SHAP & SHapley Additive exPlanations; feature attribution. \\
SMOTE & Synthetic Minority Over-sampling Technique. \\
SOFA & Sequential Organ Failure Assessment; organ dysfunction score. \\
Stay ID & Unique ICU stay identifier in MIMIC. \\
Symptom Token & Discrete representation: (symptom\_id, intensity, time, modality). \\
Threshold tuning & Choosing decision cutoff to maximize F1 or other metric. \\
Tokenization & Mapping continuous values to discrete token IDs. \\
Transfer learning & Using pretrained model on new task. \\
Validation set & Held-out data for model selection; not used for final reporting. \\
XGBoost & Extreme Gradient Boosting; tree ensemble. \\
Zero-shot & Generalization with zero labeled examples in target class. \\
\bottomrule
\end{longtable}
